\section{The Riemann integral}
For a bounded function $f$ and a partition
$P=\{a=t_0<t_1<\cdots<t_n=b\}$, let $I_j=[t_{j-1},t_j]$.
The lower and upper sums are
\[
 L(f,P)=\sum_j(\inf_{I_j}f)\abs{I_j},\qquad
 U(f,P)=\sum_j(\sup_{I_j}f)\abs{I_j},
\]
where $\abs{I_j}=t_j-t_{j-1}$ denotes interval length. The lower integral
is the supremum of all lower sums; the upper integral is the infimum of
all upper sums. Write these as $\underline I(f)$ and $\overline I(f)$.
Integrability means they are equal. A practical equivalent test is that
$U(f,P)-L(f,P)$ can be made arbitrarily small. No antiderivative formula
is assumed in the Darboux calculations below.

\exercise{32.1}{The integral of a cube from sums}
\goal Compute both Darboux integrals of $x^3$ on $[0,b]$, where $b>0$.
\step{Choose equal subintervals and identify maxima and minima.}
The cube is increasing on $[0,b]$. For $t_k=kb/n$, its maximum on
$[t_{k-1},t_k]$ is $t_k^3$ and its minimum is $t_{k-1}^3$. Therefore
\[
 U_n=\frac{b^4}{n^4}\sum_{k=1}^n k^3
     =\frac{b^4(n+1)^2}{4n^2},\qquad
 L_n=\frac{b^4}{n^4}\sum_{k=0}^{n-1} k^3
     =\frac{b^4(n-1)^2}{4n^2}.
\]
Here we used $\sum_{k=1}^n k^3=[n(n+1)/2]^2$.
\step{Squeeze the integrals.}
For every $n$,
$L_n\le\underline I(f)\le\overline I(f)\le U_n$.
Both outer sequences tend to $b^4/4$, because $(n\pm1)^2/n^2\to1$.
Thus the two integrals coincide at that value.
\answer{The function is integrable and $\int_0^b x^3\,dx=b^4/4$.}

\exercise{32.2}{A rational-only linear function}
\goal For $f(x)=x$ on rational points and $f(x)=0$ on irrational points
of $[0,b]$, compute its lower and upper integrals.
\step{Use density on each interval.}
Every nondegenerate interval $[u,v]\subseteq[0,b]$ contains irrational
points, so its infimum is zero. Its rational points can approach $v$
from below; their values approach $v$, while no value exceeds $v$.
Thus its supremum is $v$, whether or not $v$ itself is rational.
For any partition, consequently,
\[
 L(f,P)=0,\qquad U(f,P)=\sum_{j=1}^n t_j(t_j-t_{j-1}).
\]
In particular, the lower integral is $0$.
\step{Find the infimum of the upper sums.}
Since $t_j\ge t_{j-1}\ge0$,
\[
 t_j(t_j-t_{j-1})\ge
 \frac{t_j+t_{j-1}}2(t_j-t_{j-1})
 =\frac{t_j^2-t_{j-1}^2}{2}.
\]
Summing telescopes to $U(f,P)\ge b^2/2$. For the equal partition,
\[
 U(f,P_n)=\frac{b^2}{n^2}\sum_{j=1}^n j
 =\frac{b^2(n+1)}{2n}\longrightarrow\frac{b^2}{2}.
\]
An infimum bounded below by $b^2/2$ and approached by such sums equals
$b^2/2$.
\answer{(a) The lower integral is $0$ and the upper integral is $b^2/2$.
(b) They differ for $b>0$, so $f$ is not integrable.}

\exercise{32.3}{A rational-only quadratic function}
\goal Repeat the calculation for $g(x)=x^2$ on rational points and
$g(x)=0$ on irrational points of $[0,b]$.
\step{Identify each interval's extrema.}
Density and nonnegativity give $\inf_{[u,v]}g=0$ and
$\sup_{[u,v]}g=v^2$. Thus all lower sums vanish and
$U(g,P)=\sum_j t_j^2(t_j-t_{j-1})$.
\step{Compute the infimum of these upper sums.}
For $0\le u\le v$ we have $v^2\ge(v^2+vu+u^2)/3$, whence
\[
 v^2(v-u)\ge\frac{v^3-u^3}{3}.
\]
This telescopes to $U(g,P)\ge b^3/3$ for every $P$.
For equal subintervals,
\[
 U(g,P_n)=\frac{b^3}{n^3}\sum_{j=1}^n j^2
 =\frac{b^3n(n+1)(2n+1)}{6n^3}\longrightarrow\frac{b^3}{3}.
\]
\answer{(a) The lower and upper integrals are $0$ and $b^3/3$.
(b) The function is not integrable on $[0,b]$ for $b>0$.}

\exercise{32.4}{Refinement by induction}
\goal Extend the one-point refinement argument in Lemma 32.2 to any
finite refinement $P\subseteq Q$.
\step{State precisely what is being inducted on.}
Let $r$ be the number of points of $Q$ not already in $P$.
We claim
\[
 L(f,P)\le L(f,Q)\le U(f,Q)\le U(f,P).
\]
For $r=0$ the partitions are equal; the middle inequality follows
because each interval's infimum is at most its supremum.
\step{Insert one more point.}
Suppose the result is established for $r$ extra points. If $Q$ has
$r+1$ extra points, choose one of them, $q$, and set $Q'=Q\setminus\{q\}$.
By the induction hypothesis,
$L(f,P)\le L(f,Q')$ and $U(f,Q')\le U(f,P)$.
The proved one-point case, applied to $Q'$ and $Q$, gives
$L(f,Q')\le L(f,Q)$ and $U(f,Q)\le U(f,Q')$.
Combining the inequalities proves the two outer comparisons for $Q$.
The middle comparison always holds interval by interval.
\answer{Adding finitely many partition points can only increase a lower
sum and decrease an upper sum.}

\exercise{32.5}{Why the lower integral cannot exceed the upper}
\goal Apply the order property of suprema and infima to Theorem 32.4.
\step{Name the two sets.}
Let
\[
 S=\{L(f,P):P\text{ is a partition of }[a,b]\},\qquad
 T=\{U(f,Q):Q\text{ is a partition of }[a,b]\}.
\]
Both are nonempty and bounded because $f$ is bounded and the total
interval length is $b-a$.
\step{Compare sums from different partitions.}
For any $P,Q$, their union $R=P\cup Q$ refines both. Lemma 32.2 gives
\[
 L(f,P)\le L(f,R)\le U(f,R)\le U(f,Q).
\]
Thus every $s\in S$ is at most every $t\in T$. By Exercise 4.8,
$\sup S\le\inf T$. To see the order reasoning directly, each $t\in T$
is an upper bound for $S$, so $\sup S\le t$ for all $t\in T$;
therefore $\sup S$ is a lower bound for $T$, and hence is at most
$\inf T$.
\answer{$\underline I(f)=\sup S\le\inf T=\overline I(f)$.}

\exercise{32.6}{Upper and lower sums whose gap tends to zero}
\goal Prove integrability and identify the limits of the given sums.
The sums $U_n$ and $L_n$ need not come from the same partition.
Nevertheless, the cross-partition comparison just established gives
\[
 L_n\le\underline I(f)\le\overline I(f)\le U_n.
\]
Subtracting shows
$0\le\overline I(f)-\underline I(f)\le U_n-L_n\to0$.
The fixed middle difference must therefore be zero. Write the common
integral as $I$. Then
\[
 0\le U_n-I\le U_n-L_n,\qquad
 0\le I-L_n\le U_n-L_n.
\]
Each right-hand side tends to zero, so both $U_n\to I$ and $L_n\to I$.
\answer{$f$ is integrable and $\int_a^b f=\lim_n U_n=\lim_n L_n$.}

\exercise{32.7}{Changing finitely many values}
\goal Show that changing an integrable function at finitely many points
preserves its integral.
\step{First understand the finite-support error.}
Put $h=g-f$. It is zero except at finitely many points
$c_1,\ldots,c_m$, so it is bounded; choose $H>0$ with $\abs{h}\le H$.
Given any $\eta>0$, make a partition that contains the exceptional points
and nearby points on each available side, so the total length of cells
meeting these points is less than $\eta$. This is possible because there
are only finitely many points. If a point is $a$ or $b$, only its one
adjacent side is needed. On every other cell $h=0$; on the exceptional
cells $-H\le h\le H$. Thus
\[
 -H\eta\le L(h,P)\le U(h,P)\le H\eta.
\]
This proves that a finite-support function has lower and upper integrals
both zero, by letting $\eta\downarrow0$.
\step{Combine the estimate with the original integrable function.}
Here is a direct argument that does not assume a later linearity theorem.
Given $\varepsilon>0$, choose a partition for $f$ with upper-minus-lower
gap less than $\varepsilon/2$. Refine it to the partition just described,
with exceptional-cell total length less than $\eta$ and
$2H\eta<\varepsilon/2$. Refinement preserves the bound for $f$.
On any cell $J$, $g=f+h$ implies
\[
 \sup_J g\le\sup_J f+\sup_J h,\qquad
 \inf_J g\ge\inf_J f+\inf_J h.
\]
Consequently $U(g,P)-L(g,P)<\varepsilon/2+2H\eta<\varepsilon$.
This proves $g$ is integrable.
\step{Compare the two integral values.}
For the same partition,
$L(g,P)\ge L(f,P)-H\eta$ and $U(g,P)\le U(f,P)+H\eta$.
Both integrals lie between the respective sums. Hence their difference
in absolute value is at most
$U(f,P)-L(f,P)+H\eta$.
The original gap and $\eta$ can both be made arbitrarily small.
\answer{$\int_a^b g=\int_a^b f$. Isolated values do not change area.}

\exercise{32.8}{Restriction to a smaller interval}
\goal Prove integrability of $f$ on $[c,d]\subseteq[a,b]$.
Assume $c<d$; the zero-length case has integral zero by convention.
Given $\varepsilon>0$, choose a partition $P$ of $[a,b]$ with
$U(f,P)-L(f,P)<\varepsilon$. Add $c$ and $d$ to form the refinement
$Q$. Its gap is still less than $\varepsilon$.
The points of $Q$ lying in $[c,d]$ form a partition $Q_0$ of that interval.
Its gap is a sub-sum of the gap for $Q$:
\[
 U_{[c,d]}(f,Q_0)-L_{[c,d]}(f,Q_0)
 =\sum_{J\subseteq[c,d]}(\sup_J f-\inf_J f)\abs J
 \le U_{[a,b]}(f,Q)-L_{[a,b]}(f,Q).
\]
Every term omitted on the right is nonnegative. Thus the smaller
interval has partitions with gap below every positive $\varepsilon$.
\answer{The Darboux criterion proves that $f|_{[c,d]}$ is integrable.}
